% Conclusion Section - AEM Version
% Molecular engineering of semi-transparent non-fullerene acceptors for quantum-enhanced agrivoltaics

\section{Conclusion}\label{sec:Conclusion}

We have demonstrated that molecular engineering of semi-transparent non-fullerene acceptors enables quantum-enhanced energy harvesting with power conversion efficiencies reaching \SI{18.8}{\percent}. By strategically tuning the spectral transmission to target vibronic resonances in photosynthetic complexes, these materials simultaneously achieve high electrical output and preserved biological activity.

The quantum enhancement mechanism---extending coherence lifetimes by \SIrange{20}{50}{\percent} through non-Markovian dynamics---provides a new design principle for OPV materials in dual-use applications. Multi-objective optimization establishes quantitative trade-offs between photovoltaic performance and spectral transmission, with validated design rules for targeting specific applications:
\begin{itemize}
    \item Balanced configuration: \SI{18.2}{\percent} PCE with \SI{25}{\percent} ETR enhancement;
    \item Energy-focused configuration: \SI{22.1}{\percent} PCE with \SI{12}{\percent} ETR enhancement;
    \item Agriculture-focused configuration: \SI{15.4}{\percent} PCE with \SI{33}{\percent} ETR enhancement.
\end{itemize}

The AI-driven discovery pipeline screened \num{2847} candidates using equivariant graph neural networks, achieving \num{100}$\times$ acceleration compared to exhaustive DFT + HEOM simulation. This methodology provides a scalable framework for multi-objective materials optimization applicable beyond agrivoltaics.

Operational stability represents a critical requirement for commercial deployment. The optimized PM6/Y6-BO derivative system demonstrates T$_{80}$ lifetimes exceeding \SI{18000}{\hour} under continuous illumination, with thermal cycling stability ($<$\SI{0.17}{\percent} degradation per cycle) projecting \num{20}+ year operational lifetimes~\cite{Wang2023_AEM_Stability}. Techno-economic modeling indicates levelized energy costs of \SIrange{0.04}{0.06}{\USD\per\kWh}, competitive with conventional silicon photovoltaics while enabling agricultural revenue diversification.

Key structure-property relationships established in this work provide quantitative design rules for next-generation semi-transparent OPVs:
\begin{itemize}
    \item Side-chain engineering for spectral transmission tuning (\SIlist{750;820}{\nano\meter} windows);
    \item Energy level alignment optimization ($\Delta E_{\rm LUMO} = \SI{0.75}{\electronvolt}$, $\Delta E_{\rm HOMO} = \SI{0.35}{\electronvolt}$);
    \item Morphology control for balanced charge transport ($\mu_h/\mu_e \approx \num{1.2}$);
    \item Interfacial ESP gradients $>$\SI{0.3}{\electronvolt} for efficient exciton dissociation.
\end{itemize}

Geographic validation across nine climate zones, including sub-Saharan Africa, confirms persistent quantum advantages of \SIrange{18}{26}{\percent} under diverse environmental conditions. The temperature optimum of \SIrange{285}{300}{\kelvin} coincides with typical agricultural conditions, ensuring robust performance in target deployment regions.

These findings establish spectral bath engineering as a materials design paradigm for co-optimizing energy conversion and biological compatibility in next-generation semi-transparent photovoltaics. The integration of AI-driven discovery, quantum dynamics simulation, and techno-economic analysis provides a comprehensive framework applicable to artificial photosynthesis, quantum solar cells, and bio-inspired molecular electronics.
